\documentclass[12pt,a4paper]{article}

% =========================
% Font và ngôn ngữ tiếng Việt
% =========================
\usepackage{fontspec}
\usepackage{polyglossia}
\setmainlanguage{vietnamese}
\setmainfont{Times New Roman}

% =========================
% Hình ảnh và bảng
% =========================
\usepackage{graphicx}
\usepackage{longtable}
\usepackage{booktabs}
\usepackage{float}

% =========================
% Định dạng văn bản
% =========================
\usepackage{parskip}
\usepackage{setspace}
\usepackage{enumitem}

% =========================
% Hộp nội dung
% =========================
\usepackage[most]{tcolorbox}

% =========================
% TikZ
% =========================
\usepackage{tikz}
\usetikzlibrary{calc}
\usepackage{pgfornament}

% =========================
% Căn lề
% =========================
\usepackage[
    a4paper,
    top=2cm,
    bottom=2cm,
    left=2cm,
    right=2cm
]{geometry}

% =========================
% Liên kết và mục lục
% =========================
\usepackage[
    colorlinks=true,
    linkcolor=black,
    urlcolor=blue,
    citecolor=black
]{hyperref}

% =========================================================
% Định nghĩa macro \UseCaseBox dùng pgfkeys (key=value)
% =========================================================
\usepackage{pgfkeys}
\usepackage{tcolorbox}
\usepackage{enumitem}
\usepackage{caption}

\pgfkeys{
  /usecase/.is family,
  /usecase/ma/.initial={},
  /usecase/ten/.initial={},
  /usecase/tacnhan/.initial={},
  /usecase/muctieu/.initial={},
  /usecase/mota/.initial={},
  /usecase/yeucau/.initial={},
  /usecase/tiendk/.initial={},
  /usecase/dongchinh/.initial={},
  /usecase/dongphu/.initial={},
  /usecase/haudk/.initial={},
}

\NewDocumentCommand{\UseCaseBox}{m}{
  \pgfkeys{/usecase,,#1}
  \pgfkeysgetvalue{/usecase/ma}{\ucMa}
  \pgfkeysgetvalue{/usecase/ten}{\ucTen}
  \pgfkeysgetvalue{/usecase/tacnhan}{\ucTacNhan}
  \pgfkeysgetvalue{/usecase/muctieu}{\ucMucTieu}
  \pgfkeysgetvalue{/usecase/mota}{\ucMota}
  \pgfkeysgetvalue{/usecase/yeucau}{\ucYeuCau}
  \pgfkeysgetvalue{/usecase/tiendk}{\ucTienDK}
  \pgfkeysgetvalue{/usecase/dongchinh}{\ucDongChinh}
  \pgfkeysgetvalue{/usecase/dongphu}{\ucDongPhu}
  \pgfkeysgetvalue{/usecase/haudk}{\ucHauDK}

  \nopagebreak
  \begin{tcolorbox}[colback=white,colframe=black,width=\textwidth]
    \textbf{Mã ca sử dụng:} \ucMa \\
    \textbf{Tên ca sử dụng:} \ucTen \\
    \textbf{Tác nhân:} \ucTacNhan \\
    \textbf{Mục tiêu:} \ucMucTieu
    \par\vspace{0.2cm}
    \textbf{Mô tả tổng quan:} \\
    \ucMota
    \par\vspace{0.2cm}
    \textbf{Yêu cầu tham khảo:} \\
    \ucYeuCau
    \par\vspace{0.2cm}
    \textbf{Tiền điều kiện:} \\
    \ucTienDK
    \par\vspace{0.2cm}
    \textbf{Dòng sự kiện chính:}
    \ucDongChinh
    \vspace{0.2cm}
    \textbf{Các dòng sự kiện phụ:}
    \ucDongPhu
    \par\vspace{0.2cm}
    \textbf{Hậu điều kiện:} \\
    \ucHauDK
  \end{tcolorbox}
  \vspace{0.1cm}
  \captionof{table}{Đặc tả ca sử dụng \ucMa: \ucTen}\label{tab:\ucMa}
  \vspace{0.3cm}
}

\begin{document}

% ==========================================
% DANH SÁCH 36 ĐẶC TẢ CA SỬ DỤNG (USE CASE)
% ==========================================

% --- UC001 ---
\UseCaseBox{
  ma = {UC001},
  ten = {Đăng nhập},
  tacnhan = {Quản lý (Admin), Giáo viên (Teacher), Học viên (Student)},
  muctieu = {Cho phép người dùng truy cập vào hệ thống với quyền hạn tương ứng của mình.},
  mota = {Người dùng nhập tài khoản và mật khẩu; hệ thống xác thực thông tin, kiểm tra trạng thái và cấp quyền truy cập.},
  yeucau = {Yêu cầu bảo mật mật khẩu, hệ thống phân quyền.},
  tiendk = {Người dùng đã có tài khoản hợp lệ trên hệ thống.},
  dongchinh = {
    \begin{enumerate}[leftmargin=*, topsep=0pt, itemsep=0pt]
      \item Người dùng truy cập trang Đăng nhập.
      \item Người dùng nhập Tên đăng nhập/Email và Mật khẩu.
      \item Người dùng nhấn nút "Đăng nhập".
      \item Hệ thống kiểm tra thông tin đối chiếu với cơ sở dữ liệu.
      \item Hệ thống xác định vai trò và chuyển hướng người dùng đến trang chủ (Dashboard) tương ứng.
    \end{enumerate}
  },
  dongphu = {
    \begin{itemize}[leftmargin=*, topsep=0pt, itemsep=0pt]
      \item Sai thông tin: Hệ thống hiển thị lỗi "Tên đăng nhập hoặc mật khẩu không đúng".
      \item Tài khoản bị khóa: Hệ thống thông báo "Tài khoản đang bị vô hiệu hóa".
    \end{itemize}
  },
  haudk = {Phiên làm việc (session) của người dùng được tạo, người dùng có thể sử dụng các chức năng theo quyền hạn.}
}

% --- UC002 ---
\UseCaseBox{
  ma = {UC002},
  ten = {Đăng xuất},
  tacnhan = {Quản lý, Giáo viên, Học viên},
  muctieu = {Kết thúc phiên làm việc an toàn.},
  mota = {Người dùng click đăng xuất để xóa phiên đăng nhập hiện tại.},
  yeucau = {UC001 (Đăng nhập).},
  tiendk = {Người dùng đang đăng nhập vào hệ thống.},
  dongchinh = {
    \begin{enumerate}[leftmargin=*, topsep=0pt, itemsep=0pt]
      \item Người dùng chọn chức năng "Đăng xuất" trên menu.
      \item Hệ thống hủy bỏ phiên làm việc (xóa token/session).
      \item Hệ thống chuyển hướng người dùng về trang Đăng nhập.
    \end{enumerate}
  },
  dongphu = {Không có.},
  haudk = {Người dùng không thể truy cập các trang yêu cầu đăng nhập nếu không thực hiện đăng nhập lại.}
}

% --- UC003 ---
\UseCaseBox{
  ma = {UC003},
  ten = {Quản lý tài khoản cá nhân},
  tacnhan = {Quản lý, Giáo viên, Học viên},
  muctieu = {Cập nhật thông tin hồ sơ và bảo mật cá nhân.},
  mota = {Xem, chỉnh sửa thông tin cá nhân (tên, SĐT, avatar) và đổi mật khẩu.},
  yeucau = {Không.},
  tiendk = {Người dùng đã đăng nhập.},
  dongchinh = {
    \begin{enumerate}[leftmargin=*, topsep=0pt, itemsep=0pt]
      \item Người dùng truy cập trang "Hồ sơ cá nhân".
      \item Hệ thống hiển thị thông tin hiện tại.
      \item Người dùng nhập các thay đổi (thông tin hoặc mật khẩu mới).
      \item Người dùng nhấn "Lưu thay đổi".
      \item Hệ thống cập nhật DB và thông báo thành công.
    \end{enumerate}
  },
  dongphu = {
    \begin{itemize}[leftmargin=*, topsep=0pt, itemsep=0pt]
      \item Mật khẩu cũ không đúng: Yêu cầu nhập lại mật khẩu hiện tại.
      \item Dữ liệu không hợp lệ: Báo lỗi format (VD: sai định dạng số điện thoại).
    \end{itemize}
  },
  haudk = {Thông tin người dùng được cập nhật mới nhất.}
}

% --- UC004 ---
\UseCaseBox{
  ma = {UC004},
  ten = {Quản lý tài khoản người dùng},
  tacnhan = {Quản lý (Admin)},
  muctieu = {Duy trì và kiểm soát danh sách tài khoản trong hệ thống.},
  mota = {Admin thêm mới, chỉnh sửa, khóa/mở khóa tài khoản.},
  yeucau = {UC005 (Phân quyền).},
  tiendk = {Đăng nhập với quyền Quản lý.},
  dongchinh = {
    \begin{enumerate}[leftmargin=*, topsep=0pt, itemsep=0pt]
      \item Quản lý chọn menu "Quản lý tài khoản".
      \item Hệ thống hiển thị danh sách tài khoản.
      \item Quản lý chọn Thêm/Sửa/Khóa tài khoản.
      \item Quản lý nhập thông tin và xác nhận.
      \item Hệ thống lưu thay đổi và cập nhật danh sách.
    \end{enumerate}
  },
  dongphu = {
    \begin{itemize}[leftmargin=*, topsep=0pt, itemsep=0pt]
      \item Trùng email: Hệ thống báo lỗi "Email đã tồn tại".
    \end{itemize}
  },
  haudk = {Danh sách tài khoản và trạng thái được cập nhật.}
}

% --- UC005 ---
\UseCaseBox{
  ma = {UC005},
  ten = {Phân quyền người dùng},
  tacnhan = {Quản lý},
  muctieu = {Gán đúng vai trò (Role) cho tài khoản để đảm bảo tính bảo mật.},
  mota = {Cài đặt quyền (Admin, Teacher, Student) cho các tài khoản.},
  yeucau = {UC004.},
  tiendk = {Đăng nhập quyền Quản lý, tài khoản cần phân quyền đã tồn tại.},
  dongchinh = {
    \begin{enumerate}[leftmargin=*, topsep=0pt, itemsep=0pt]
      \item Quản lý chọn tài khoản cần phân quyền.
      \item Chọn chức năng "Chỉnh sửa vai trò".
      \item Chọn vai trò từ danh sách (Giáo viên, Học viên...).
      \item Nhấn "Lưu". Hệ thống cập nhật quyền cho tài khoản.
    \end{enumerate}
  },
  dongphu = {Không có.},
  haudk = {Tài khoản sở hữu các chức năng tương ứng với quyền mới gán.}
}

% --- UC006 ---
\UseCaseBox{
  ma = {UC006},
  ten = {Quản lý giáo viên},
  tacnhan = {Quản lý},
  muctieu = {Quản lý hồ sơ chuyên môn và phân công giảng dạy.},
  mota = {Xem, thêm, cập nhật hồ sơ giáo viên và trạng thái làm việc.},
  yeucau = {UC004, UC008.},
  tiendk = {Đăng nhập quyền Quản lý.},
  dongchinh = {
    \begin{enumerate}[leftmargin=*, topsep=0pt, itemsep=0pt]
      \item Truy cập "Quản lý giáo viên".
      \item Xem danh sách và thông tin (trình độ, chứng chỉ).
      \item Thêm mới hoặc cập nhật thông tin giáo viên.
      \item Lưu dữ liệu.
    \end{enumerate}
  },
  dongphu = {
    \begin{itemize}[leftmargin=*, topsep=0pt, itemsep=0pt]
      \item Hủy thao tác: Quay về màn hình trước.
    \end{itemize}
  },
  haudk = {Danh sách hồ sơ giáo viên được đồng bộ.}
}

% --- UC007 ---
\UseCaseBox{
  ma = {UC007},
  ten = {Quản lý học viên},
  tacnhan = {Quản lý},
  muctieu = {Kiểm soát thông tin đầu vào và trạng thái học tập của học viên.},
  mota = {Quản lý hồ sơ cá nhân, trình độ đầu vào của học viên.},
  yeucau = {UC004, UC009.},
  tiendk = {Đăng nhập quyền Quản lý.},
  dongchinh = {
    \begin{enumerate}[leftmargin=*, topsep=0pt, itemsep=0pt]
      \item Truy cập "Quản lý học viên".
      \item Thêm học viên đơn lẻ hoặc import từ file Excel.
      \item Cập nhật thông tin (Tên, tuổi, số điện thoại phụ huynh).
      \item Hệ thống lưu và tạo hồ sơ.
    \end{enumerate}
  },
  dongphu = {
    \begin{itemize}[leftmargin=*, topsep=0pt, itemsep=0pt]
      \item Import lỗi: File Excel sai format, hệ thống báo lỗi các dòng cụ thể.
    \end{itemize}
  },
  haudk = {Học viên có hồ sơ lưu trữ trên hệ thống.}
}

% --- UC008 ---
\UseCaseBox{
  ma = {UC008},
  ten = {Quản lý lớp học},
  tacnhan = {Quản lý},
  muctieu = {Tổ chức các lớp học và phân công giáo viên.},
  mota = {Tạo lớp, đóng/mở lớp, gán giáo viên phụ trách.},
  yeucau = {UC006.},
  tiendk = {Đăng nhập quyền Quản lý.},
  dongchinh = {
    \begin{enumerate}[leftmargin=*, topsep=0pt, itemsep=0pt]
      \item Chọn "Quản lý lớp học" $\rightarrow$ "Tạo mới".
      \item Nhập Tên lớp, Trình độ, Thời gian học.
      \item Chọn giáo viên phụ trách từ danh sách.
      \item Nhấn "Lưu".
    \end{enumerate}
  },
  dongphu = {
    \begin{itemize}[leftmargin=*, topsep=0pt, itemsep=0pt]
      \item Đóng lớp học khi đã hoàn thành khóa học.
    \end{itemize}
  },
  haudk = {Lớp học được tạo, sẵn sàng để thêm học viên.}
}

% --- UC009 ---
\UseCaseBox{
  ma = {UC009},
  ten = {Quản lý thành viên lớp học},
  tacnhan = {Quản lý, Giáo viên},
  muctieu = {Sắp xếp học viên vào đúng lớp.},
  mota = {Thêm, xóa, hoặc chuyển lớp cho học viên.},
  yeucau = {UC007, UC008.},
  tiendk = {Lớp học và hồ sơ học viên đã tồn tại.},
  dongchinh = {
    \begin{enumerate}[leftmargin=*, topsep=0pt, itemsep=0pt]
      \item Truy cập chi tiết Lớp học $\rightarrow$ "Thành viên".
      \item Nhấn "Thêm học viên". Chọn từ danh sách có sẵn.
      \item Nhấn xác nhận.
      \item Hệ thống thêm học viên vào lớp.
    \end{enumerate}
  },
  dongphu = {
    \begin{itemize}[leftmargin=*, topsep=0pt, itemsep=0pt]
      \item Chuyển học viên từ lớp A sang lớp B.
    \end{itemize}
  },
  haudk = {Danh sách (sĩ số) của lớp học được cập nhật.}
}

% --- UC010 ---
\UseCaseBox{
  ma = {UC010},
  ten = {Xem danh sách lớp học},
  tacnhan = {Giáo viên, Học viên},
  muctieu = {Biết được các lớp mình đang giảng dạy hoặc tham gia.},
  mota = {Hiển thị danh sách các lớp học được phân công.},
  yeucau = {UC008, UC009.},
  tiendk = {Đăng nhập quyền Teacher hoặc Student.},
  dongchinh = {
    \begin{enumerate}[leftmargin=*, topsep=0pt, itemsep=0pt]
      \item Truy cập menu "Lớp học của tôi".
      \item Hệ thống truy xuất DB và hiển thị danh sách lớp (Tên lớp, sĩ số, lịch học).
    \end{enumerate}
  },
  dongphu = {
    \begin{itemize}[leftmargin=*, topsep=0pt, itemsep=0pt]
      \item Nếu chưa có lớp, hệ thống thông báo "Bạn chưa tham gia lớp học nào".
    \end{itemize}
  },
  haudk = {Người dùng thấy được lớp học tương ứng.}
}

% --- UC011 ---
\UseCaseBox{
  ma = {UC011},
  ten = {Quản lý bài tập},
  tacnhan = {Giáo viên},
  muctieu = {Kiểm soát thư viện bài tập của lớp.},
  mota = {Xem, tìm kiếm, chỉnh sửa, xóa bài tập đã tạo.},
  yeucau = {UC012.},
  tiendk = {Đăng nhập quyền Teacher.},
  dongchinh = {
    \begin{enumerate}[leftmargin=*, topsep=0pt, itemsep=0pt]
      \item Vào mục "Bài tập" của lớp.
      \item Hệ thống hiển thị danh sách bài tập.
      \item Sử dụng bộ lọc (kỹ năng, trạng thái) để tìm kiếm.
      \item Chọn Sửa hoặc Xóa bài tập.
    \end{enumerate}
  },
  dongphu = {
    \begin{itemize}[leftmargin=*, topsep=0pt, itemsep=0pt]
      \item Sao chép bài tập sang lớp khác.
    \end{itemize}
  },
  haudk = {Thư viện bài tập được cập nhật/sắp xếp.}
}

% --- UC012 ---
\UseCaseBox{
  ma = {UC012},
  ten = {Tạo bài tập},
  tacnhan = {Giáo viên},
  muctieu = {Xây dựng nội dung bài kiểm tra/luyện tập cho học viên.},
  mota = {Thiết lập câu hỏi, đáp án, chọn kỹ năng, thời gian làm bài, điểm số.},
  yeucau = {UC011.},
  tiendk = {Đăng nhập quyền Teacher.},
  dongchinh = {
    \begin{enumerate}[leftmargin=*, topsep=0pt, itemsep=0pt]
      \item Chọn "Tạo bài tập mới".
      \item Nhập Tiêu đề, Chọn loại kỹ năng (Nghe/Nói/Đọc/Viết).
      \item Soạn thảo nội dung (tải file Audio, tạo câu hỏi trắc nghiệm, tạo đề tự luận).
      \item Nhập đáp án chuẩn và thiết lập điểm.
      \item Lưu lại dưới dạng bản nháp hoặc xuất bản luôn.
    \end{enumerate}
  },
  dongphu = {
    \begin{itemize}[leftmargin=*, topsep=0pt, itemsep=0pt]
      \item Import ngân hàng câu hỏi từ file Excel/Word.
    \end{itemize}
  },
  haudk = {Bài tập được lưu vào cơ sở dữ liệu hệ thống.}
}

% --- UC013 ---
\UseCaseBox{
  ma = {UC013},
  ten = {Mở bài tập},
  tacnhan = {Giáo viên},
  muctieu = {Giao bài cho học viên bắt đầu làm.},
  mota = {Đổi trạng thái bài tập sang "Đang mở", thiết lập deadline.},
  yeucau = {UC012, UC015.},
  tiendk = {Bài tập đã được tạo xong nội dung.},
  dongchinh = {
    \begin{enumerate}[leftmargin=*, topsep=0pt, itemsep=0pt]
      \item Chọn bài tập trong danh sách.
      \item Chọn "Giao bài / Mở bài".
      \item Thiết lập thời gian bắt đầu và hạn chót (Deadline).
      \item Xác nhận. Hệ thống gửi thông báo cho học viên.
    \end{enumerate}
  },
  dongphu = {
    \begin{itemize}[leftmargin=*, topsep=0pt, itemsep=0pt]
      \item Hẹn giờ mở bài tự động.
    \end{itemize}
  },
  haudk = {Học viên có thể nhìn thấy và truy cập làm bài tập.}
}

% --- UC014 ---
\UseCaseBox{
  ma = {UC014},
  ten = {Khóa bài tập},
  tacnhan = {Giáo viên (hoặc hệ thống tự động)},
  muctieu = {Dừng việc nhận bài làm mới từ học viên.},
  mota = {Vô hiệu hóa chức năng nộp bài khi hết hạn hoặc khi giáo viên chủ động khóa.},
  yeucau = {UC013.},
  tiendk = {Bài tập đang ở trạng thái Mở.},
  dongchinh = {
    \begin{enumerate}[leftmargin=*, topsep=0pt, itemsep=0pt]
      \item Giáo viên chọn bài tập đang mở.
      \item Chọn chức năng "Khóa bài tập".
      \item Hệ thống chặn nút nộp bài đối với học viên.
    \end{enumerate}
  },
  dongphu = {
    \begin{itemize}[leftmargin=*, topsep=0pt, itemsep=0pt]
      \item Hệ thống tự động khóa khi tới thời điểm Deadline đã thiết lập.
    \end{itemize}
  },
  haudk = {Bài tập chuyển sang trạng thái "Đã đóng", học viên không thể nộp bài.}
}

% --- UC015 ---
\UseCaseBox{
  ma = {UC015},
  ten = {Xem bài tập},
  tacnhan = {Học viên},
  muctieu = {Nắm bắt yêu cầu và deadline của bài tập được giao.},
  mota = {Học viên xem danh sách bài tập, trạng thái và yêu cầu chi tiết.},
  yeucau = {UC013.},
  tiendk = {Đăng nhập quyền Student.},
  dongchinh = {
    \begin{enumerate}[leftmargin=*, topsep=0pt, itemsep=0pt]
      \item Học viên truy cập "Bài tập của tôi".
      \item Hệ thống hiển thị danh sách (Tiêu đề, Deadline, Trạng thái: Chưa làm, Đã nộp...).
      \item Nhấn vào một bài tập để xem hướng dẫn chi tiết.
    \end{enumerate}
  },
  dongphu = {
    \begin{itemize}[leftmargin=*, topsep=0pt, itemsep=0pt]
      \item Lọc bài tập theo kỹ năng (Nghe, Nói, Đọc, Viết).
    \end{itemize}
  },
  haudk = {Học viên hiểu yêu cầu và sẵn sàng làm bài.}
}

% --- UC016 ---
\UseCaseBox{
  ma = {UC016},
  ten = {Làm bài tập viết},
  tacnhan = {Học viên},
  muctieu = {Hoàn thành phần trả lời cho kỹ năng Viết.},
  mota = {Học viên đọc đề bài, soạn thảo nội dung trực tiếp trên web hoặc tải file lên.},
  yeucau = {UC015, UC020, UC021.},
  tiendk = {Bài tập Viết đang mở.},
  dongchinh = {
    \begin{enumerate}[leftmargin=*, topsep=0pt, itemsep=0pt]
      \item Học viên mở bài tập Viết.
      \item Đọc đề bài (Task 1, Task 2).
      \item Nhập văn bản vào khung soạn thảo (có bộ đếm từ - word count).
      \item Kiểm tra lại văn bản.
    \end{enumerate}
  },
  dongphu = {
    \begin{itemize}[leftmargin=*, topsep=0pt, itemsep=0pt]
      \item Upload file (.docx, .pdf) thay vì gõ trực tiếp.
    \end{itemize}
  },
  haudk = {Nội dung viết được điền sẵn sàng để lưu nháp hoặc nộp.}
}

% --- UC017 ---
\UseCaseBox{
  ma = {UC017},
  ten = {Làm bài tập nói},
  tacnhan = {Học viên},
  muctieu = {Hoàn thành bài kiểm tra kỹ năng Nói.},
  mota = {Ghi âm trực tiếp qua trình duyệt hoặc tải file âm thanh lên hệ thống.},
  yeucau = {UC015, UC021.},
  tiendk = {Bài tập Nói đang mở, thiết bị có Microphone hợp lệ.},
  dongchinh = {
    \begin{enumerate}[leftmargin=*, topsep=0pt, itemsep=0pt]
      \item Mở bài tập Nói, đọc/nghe câu hỏi.
      \item Nhấn "Ghi âm" (Cấp quyền Micro cho trình duyệt).
      \item Trả lời câu hỏi.
      \item Nhấn "Dừng" và nghe lại đoạn ghi âm.
    \end{enumerate}
  },
  dongphu = {
    \begin{itemize}[leftmargin=*, topsep=0pt, itemsep=0pt]
      \item Tải lên tệp âm thanh có sẵn (.mp3, .wav).
      \item Ghi âm lại nhiều lần trước khi nộp.
    \end{itemize}
  },
  haudk = {File âm thanh được tải lên máy chủ chờ nộp bài.}
}

% --- UC018 ---
\UseCaseBox{
  ma = {UC018},
  ten = {Làm bài tập đọc},
  tacnhan = {Học viên},
  muctieu = {Kiểm tra kỹ năng đọc hiểu.},
  mota = {Đọc văn bản, trả lời các dạng câu hỏi trắc nghiệm, điền từ, True/False/Not Given.},
  yeucau = {UC015, UC021.},
  tiendk = {Bài tập Đọc đang mở.},
  dongchinh = {
    \begin{enumerate}[leftmargin=*, topsep=0pt, itemsep=0pt]
      \item Mở bài tập Đọc.
      \item Đọc đoạn văn hiển thị bên trái màn hình.
      \item Trả lời các câu hỏi tương ứng bên phải màn hình.
    \end{enumerate}
  },
  dongphu = {
    \begin{itemize}[leftmargin=*, topsep=0pt, itemsep=0pt]
      \item Highlight (bôi màu) các đoạn văn bản nháp trực tiếp trên trình duyệt.
    \end{itemize}
  },
  haudk = {Các câu trả lời được chọn đầy đủ.}
}

% --- UC019 ---
\UseCaseBox{
  ma = {UC019},
  ten = {Làm bài tập nghe},
  tacnhan = {Học viên},
  muctieu = {Kiểm tra kỹ năng nghe hiểu âm thanh tiếng Anh.},
  mota = {Học viên phát file audio và chọn/nhập đáp án cùng lúc.},
  yeucau = {UC015, UC021.},
  tiendk = {Bài tập Nghe đang mở, thiết bị có loa/tai nghe hoạt động.},
  dongchinh = {
    \begin{enumerate}[leftmargin=*, topsep=0pt, itemsep=0pt]
      \item Mở bài tập Nghe.
      \item Nhấn nút "Play" (Phát) file Audio.
      \item Nghe và chọn/điền đáp án vào các câu hỏi bên dưới.
    \end{enumerate}
  },
  dongphu = {
    \begin{itemize}[leftmargin=*, topsep=0pt, itemsep=0pt]
      \item Giáo viên giới hạn số lần Play audio (vd: chỉ được nghe 1 lần) $\rightarrow$ Hệ thống tự động disable nút Play sau khi hết lượt.
    \end{itemize}
  },
  haudk = {Câu trả lời được nhập đầy đủ, kết thúc thời gian nghe.}
}

% --- UC020 ---
\UseCaseBox{
  ma = {UC020},
  ten = {Lưu và khôi phục bài làm},
  tacnhan = {Học viên},
  muctieu = {Bảo vệ dữ liệu bài làm khỏi sự cố (mất mạng, tắt máy).},
  mota = {Hệ thống tự động auto-save bài làm sau mỗi vài giây, hoặc lưu thủ công.},
  yeucau = {UC016-UC019.},
  tiendk = {Đang trong quá trình làm bài.},
  dongchinh = {
    \begin{enumerate}[leftmargin=*, topsep=0pt, itemsep=0pt]
      \item Học viên đang gõ/nhập đáp án.
      \item Hệ thống chạy ngầm chức năng tự động lưu (Auto-save) lên local storage / database sau mỗi 30s.
      \item Nếu học viên đóng trình duyệt và mở lại, hệ thống gợi ý "Khôi phục bài làm nháp gần nhất".
      \item Học viên chọn "Đồng ý", dữ liệu cũ được điền lại.
    \end{enumerate}
  },
  dongphu = {
    \begin{itemize}[leftmargin=*, topsep=0pt, itemsep=0pt]
      \item Học viên tự nhấn nút "Lưu nháp" thủ công.
    \end{itemize}
  },
  haudk = {Dữ liệu không bị mất trong quá trình thực hiện bài tập.}
}

% --- UC021 ---
\UseCaseBox{
  ma = {UC021},
  ten = {Nộp bài tập},
  tacnhan = {Học viên},
  muctieu = {Gửi bài làm lên hệ thống để chờ chấm điểm.},
  mota = {Học viên xác nhận hoàn thành và nộp; hệ thống ghi nhận thời gian nộp.},
  yeucau = {UC016-UC019, UC023.},
  tiendk = {Đã hoàn thành các câu hỏi, bài tập chưa bị khóa.},
  dongchinh = {
    \begin{enumerate}[leftmargin=*, topsep=0pt, itemsep=0pt]
      \item Học viên nhấn nút "Nộp bài".
      \item Hệ thống hiển thị popup xác nhận "Bạn có chắc chắn muốn nộp bài?".
      \item Học viên nhấn "Xác nhận".
      \item Hệ thống ghi nhận trạng thái bài nộp, thời gian nộp và khóa bài làm của cá nhân học viên.
    \end{enumerate}
  },
  dongphu = {
    \begin{itemize}[leftmargin=*, topsep=0pt, itemsep=0pt]
      \item Chưa làm hết câu hỏi $\rightarrow$ Hệ thống cảnh báo "Bạn còn X câu chưa làm, vẫn tiếp tục nộp?".
    \end{itemize}
  },
  haudk = {Bài làm được gửi thành công, chuyển sang trạng thái "Chờ chấm".}
}

% --- UC022 ---
\UseCaseBox{
  ma = {UC022},
  ten = {Quản lý lần làm và nộp lại bài},
  tacnhan = {Học viên, Giáo viên},
  muctieu = {Cho phép làm lại bài nếu chính sách bài tập cho phép.},
  mota = {Hệ thống kiểm tra số lần được phép làm; học viên có thể nộp lại.},
  yeucau = {UC021.},
  tiendk = {Bài tập được thiết lập số lần làm $>$ 1 và chưa quá hạn.},
  dongchinh = {
    \begin{enumerate}[leftmargin=*, topsep=0pt, itemsep=0pt]
      \item Học viên đã nộp bài lần 1.
      \item Học viên vào lại bài tập, thấy nút "Làm lại bài (Còn x lượt)".
      \item Chọn "Làm lại", hệ thống tạo form bài làm mới trắng trơn.
      \item Thực hiện và nộp bài lần 2.
    \end{enumerate}
  },
  dongphu = {
    \begin{itemize}[leftmargin=*, topsep=0pt, itemsep=0pt]
      \item Hết số lượt làm lại $\rightarrow$ Nút bị ẩn.
    \end{itemize}
  },
  haudk = {Lưu trữ lịch sử tất cả các lần nộp của học viên.}
}

% --- UC023 ---
\UseCaseBox{
  ma = {UC023},
  ten = {Xem trạng thái bài làm},
  tacnhan = {Học viên},
  muctieu = {Theo dõi tiến trình bài tập của mình.},
  mota = {Cập nhật hiển thị (Chưa làm, Đã nộp, Đang chấm, Đã chấm).},
  yeucau = {UC021, UC024.},
  tiendk = {Đã được giao bài tập.},
  dongchinh = {
    \begin{enumerate}[leftmargin=*, topsep=0pt, itemsep=0pt]
      \item Học viên truy cập trang tổng quan (Dashboard) hoặc danh sách bài tập.
      \item Hệ thống hiển thị nhãn trạng thái (Status tag) màu sắc bên cạnh mỗi bài tập.
    \end{enumerate}
  },
  dongphu = {Không có.},
  haudk = {Học viên biết rõ bài nào cần làm, bài nào đã có kết quả.}
}

% --- UC024 ---
\UseCaseBox{
  ma = {UC024},
  ten = {Chấm bài},
  tacnhan = {Giáo viên},
  muctieu = {Đánh giá năng lực của học viên và trả kết quả.},
  mota = {Giáo viên xem bài làm, chấm điểm, nhận xét thủ công.},
  yeucau = {UC021.},
  tiendk = {Đăng nhập Teacher, Học viên đã nộp bài.},
  dongchinh = {
    \begin{enumerate}[leftmargin=*, topsep=0pt, itemsep=0pt]
      \item Giáo viên chọn danh sách bài "Chờ chấm".
      \item Mở bài làm của 1 học viên.
      \item Đọc/nghe bài làm, nhập điểm cho từng câu/phần.
      \item Ghi chú, comment nhận xét vào trực tiếp bài làm (inline comment) hoặc nhận xét chung.
      \item Nhấn "Xác nhận \& Trả bài".
    \end{enumerate}
  },
  dongphu = {
    \begin{itemize}[leftmargin=*, topsep=0pt, itemsep=0pt]
      \item Lưu nháp kết quả chấm để xem lại sau.
    \end{itemize}
  },
  haudk = {Trạng thái bài làm đổi thành "Đã chấm", học viên nhận được thông báo có điểm.}
}

% --- UC025 ---
\UseCaseBox{
  ma = {UC025},
  ten = {Hỗ trợ chấm chữa bài bằng AI (Chung)},
  tacnhan = {Giáo viên, Hệ thống AI},
  muctieu = {Tối ưu hóa thời gian chấm bài cho giáo viên, đưa ra đánh giá khách quan.},
  mota = {Gọi API AI để phân tích bài làm, đề xuất lỗi và điểm, sau đó giáo viên duyệt lại.},
  yeucau = {UC026, UC027.},
  tiendk = {Hệ thống được tích hợp module AI hoạt động ổn định.},
  dongchinh = {
    \begin{enumerate}[leftmargin=*, topsep=0pt, itemsep=0pt]
      \item Giáo viên mở bài làm của học viên.
      \item Nhấn "Gợi ý chấm bằng AI".
      \item Hệ thống gửi dữ liệu lên AI model, chờ trả kết quả.
      \item Hệ thống hiển thị điểm đề xuất, các lỗi highlight trên giao diện.
      \item Giáo viên kiểm tra, có thể sửa lại điểm AI đề xuất nếu thấy chưa phù hợp.
      \item Xác nhận kết quả cuối.
    \end{enumerate}
  },
  dongphu = {
    \begin{itemize}[leftmargin=*, topsep=0pt, itemsep=0pt]
      \item Lỗi kết nối AI $\rightarrow$ Báo lỗi "Hệ thống AI đang gián đoạn, vui lòng chấm thủ công".
    \end{itemize}
  },
  haudk = {Bài làm có kết quả đánh giá sơ bộ từ AI được lưu lại.}
}

% --- UC026 ---
\UseCaseBox{
  ma = {UC026},
  ten = {Hỗ trợ chấm chữa bài viết bằng AI},
  tacnhan = {Giáo viên, AI},
  muctieu = {Phát hiện lỗi ngữ pháp, từ vựng và cấu trúc bài viết tự động.},
  mota = {AI đánh giá Writing theo các tiêu chí (Grammar, Lexical Resource, Coherence, Task Achievement).},
  yeucau = {UC025.},
  tiendk = {Bài làm thuộc kỹ năng Viết.},
  dongchinh = {
    \begin{enumerate}[leftmargin=*, topsep=0pt, itemsep=0pt]
      \item Gửi text bài viết lên AI.
      \item AI trả về kết quả bôi đỏ (lỗi chính tả), bôi vàng (lỗi ngữ pháp), bôi xanh lam (từ vựng chưa hay).
      \item AI đề xuất cách sửa (Rewrite suggestion) cho từng câu lỗi.
      \item AI tổng hợp điểm từng tiêu chí.
      \item Giáo viên duyệt, chọn Accept (chấp nhận) hoặc Reject (từ chối) gợi ý của AI.
    \end{enumerate}
  },
  dongphu = {
    \begin{itemize}[leftmargin=*, topsep=0pt, itemsep=0pt]
      \item Bài viết phát hiện đạo văn (Plagiarism) $\rightarrow$ AI cảnh báo đỏ 100\%.
    \end{itemize}
  },
  haudk = {Bài viết có đầy đủ chú thích lỗi và điểm chi tiết.}
}

% --- UC027 ---
\UseCaseBox{
  ma = {UC027},
  ten = {Hỗ trợ chấm chữa bài nói bằng AI},
  tacnhan = {Giáo viên, AI},
  muctieu = {Đánh giá phát âm, độ lưu loát của giọng nói tự động.},
  mota = {AI chuyển Speech-to-text và phân tích phát âm (Pronunciation), độ trôi chảy (Fluency).},
  yeucau = {UC025.},
  tiendk = {Bài làm thuộc kỹ năng Nói, chất lượng âm thanh rõ.},
  dongchinh = {
    \begin{enumerate}[leftmargin=*, topsep=0pt, itemsep=0pt]
      \item Hệ thống gửi file Audio lên AI.
      \item AI bóc băng (Transcript) đoạn nói của học viên.
      \item Phân tích các từ phát âm sai (mispronounced words) hiển thị màu đỏ trên Transcript.
      \item Đánh giá nhịp điệu (Intonation) và độ ngập ngừng. Đề xuất điểm.
      \item Giáo viên nghe lại và chốt điểm cuối.
    \end{enumerate}
  },
  dongphu = {
    \begin{itemize}[leftmargin=*, topsep=0pt, itemsep=0pt]
      \item Tạp âm quá ồn $\rightarrow$ AI trả về cảnh báo "Chất lượng âm thanh quá thấp không thể phân tích".
    \end{itemize}
  },
  haudk = {Đánh giá kỹ năng Speaking hoàn tất.}
}

% --- UC028 ---
\UseCaseBox{
  ma = {UC028},
  ten = {Chấm bài đọc (Reading)},
  tacnhan = {Hệ thống, Giáo viên},
  muctieu = {Tự động chấm điểm khách quan.},
  mota = {Hệ thống so khớp đáp án của học viên với đáp án gốc (Answer key) của giáo viên.},
  yeucau = {UC018.},
  tiendk = {Đã khai báo đáp án đúng khi tạo bài.},
  dongchinh = {
    \begin{enumerate}[leftmargin=*, topsep=0pt, itemsep=0pt]
      \item Học viên nộp bài Đọc.
      \item Hệ thống tự động so sánh từng câu trả lời.
      \item Tính tổng số câu đúng và quy đổi ra điểm hệ 10 hoặc hệ IELTS/TOEIC.
      \item Lưu điểm lập tức vào hệ thống.
    \end{enumerate}
  },
  dongphu = {
    \begin{itemize}[leftmargin=*, topsep=0pt, itemsep=0pt]
      \item Có câu trả lời dạng điền từ (tương đối đúng) $\rightarrow$ Hệ thống gán trạng thái "Cần giáo viên duyệt lại".
    \end{itemize}
  },
  haudk = {Điểm số phần Reading được xác định.}
}

% --- UC029 ---
\UseCaseBox{
  ma = {UC029},
  ten = {Chấm bài nghe (Listening)},
  tacnhan = {Hệ thống, Giáo viên},
  muctieu = {Tự động hóa chấm điểm nghe.},
  mota = {Tương tự bài đọc, hệ thống đối chiếu và tính điểm tự động.},
  yeucau = {UC019.},
  tiendk = {Khai báo sẵn Answer key.},
  dongchinh = {
    \begin{enumerate}[leftmargin=*, topsep=0pt, itemsep=0pt]
      \item Học viên nộp bài Nghe.
      \item Hệ thống tự động chấm câu trắc nghiệm, điền từ (khớp chính xác).
      \item Đưa ra điểm tự động.
    \end{enumerate}
  },
  dongphu = {
    \begin{itemize}[leftmargin=*, topsep=0pt, itemsep=0pt]
      \item Nếu học viên điền sai viết hoa/thường nhưng đúng từ $\rightarrow$ Cấu hình cho phép tính là Đúng.
    \end{itemize}
  },
  haudk = {Điểm Listening được ghi nhận.}
}

% --- UC030 ---
\UseCaseBox{
  ma = {UC030},
  ten = {Xem điểm số},
  tacnhan = {Học viên, Giáo viên},
  muctieu = {Nắm bắt điểm số của các bài tập đã nộp.},
  mota = {Hiển thị điểm thi/điểm bài tập trực quan.},
  yeucau = {Không.},
  tiendk = {Bài đã được chấm và công bố điểm.},
  dongchinh = {
    \begin{enumerate}[leftmargin=*, topsep=0pt, itemsep=0pt]
      \item Truy cập mục "Bảng điểm / Kết quả".
      \item Hệ thống hiển thị danh sách bài tập kèm điểm số đạt được, điểm tối đa.
    \end{enumerate}
  },
  dongphu = {Không.},
  haudk = {Thông tin điểm số minh bạch đến người dùng.}
}

% --- UC031 ---
\UseCaseBox{
  ma = {UC031},
  ten = {Quản lý điểm số},
  tacnhan = {Giáo viên},
  muctieu = {Quản lý sổ điểm của cả lớp.},
  mota = {Xem bảng điểm tổng hợp, cập nhật, chỉnh sửa điểm nếu có sai sót.},
  yeucau = {UC030.},
  tiendk = {Đăng nhập Teacher.},
  dongchinh = {
    \begin{enumerate}[leftmargin=*, topsep=0pt, itemsep=0pt]
      \item Vào mục "Quản lý điểm số" của 1 lớp.
      \item Hệ thống hiển thị bảng Matrix (Học viên x Bài tập).
      \item Giáo viên có thể click vào ô điểm để sửa thủ công (kèm lý do sửa).
      \item Lưu thay đổi và xuất file Excel nếu cần.
    \end{enumerate}
  },
  dongphu = {Không.},
  haudk = {Sổ điểm chuẩn xác được lưu vết (audit log).}
}

% --- UC032 ---
\UseCaseBox{
  ma = {UC032},
  ten = {Xem kết quả học tập},
  tacnhan = {Học viên},
  muctieu = {Nhìn nhận bức tranh toàn cảnh về tiến bộ cá nhân.},
  mota = {Xem biểu đồ, điểm trung bình, xếp hạng trong lớp.},
  yeucau = {UC030.},
  tiendk = {Có dữ liệu điểm số hệ thống.},
  dongchinh = {
    \begin{enumerate}[leftmargin=*, topsep=0pt, itemsep=0pt]
      \item Học viên chọn "Tiến độ học tập".
      \item Hệ thống render biểu đồ radar (4 kỹ năng Nghe-Nói-Đọc-Viết) thể hiện điểm yếu/mạnh.
      \item Xem điểm trung bình toàn khóa.
    \end{enumerate}
  },
  dongphu = {Không.},
  haudk = {Học viên có định hướng cải thiện việc học.}
}

% --- UC033 ---
\UseCaseBox{
  ma = {UC033},
  ten = {Xem bài chữa},
  tacnhan = {Học viên},
  muctieu = {Học hỏi từ các lỗi sai để rút kinh nghiệm.},
  mota = {Học viên xem lại chi tiết bài làm, đọc comment nhận xét của giáo viên và AI.},
  yeucau = {UC024, UC025.},
  tiendk = {Bài đã được chấm.},
  dongchinh = {
    \begin{enumerate}[leftmargin=*, topsep=0pt, itemsep=0pt]
      \item Học viên click vào "Chi tiết bài chữa" của bài đã có điểm.
      \item Hệ thống hiển thị bài làm kèm các đánh dấu màu sắc (highlight lỗi).
      \item Đọc/nghe phần nhận xét (Feedback) từ giáo viên.
      \item Xem đáp án chuẩn (Answer key/Sample answer).
    \end{enumerate}
  },
  dongphu = {Không.},
  haudk = {Học viên tiếp thu kiến thức sửa lỗi.}
}

% --- UC034 ---
\UseCaseBox{
  ma = {UC034},
  ten = {Xem và đánh giá kết quả học viên},
  tacnhan = {Giáo viên},
  muctieu = {Đánh giá định kỳ năng lực từng cá nhân.},
  mota = {Giáo viên xem hồ sơ tiến bộ của một học viên cụ thể và viết nhận xét cuối tháng/khóa.},
  yeucau = {UC032.},
  tiendk = {Đăng nhập Teacher.},
  dongchinh = {
    \begin{enumerate}[leftmargin=*, topsep=0pt, itemsep=0pt]
      \item Chọn một học viên trong danh sách lớp.
      \item Xem lịch sử làm bài, biểu đồ điểm số các kỹ năng của học viên đó.
      \item Viết đánh giá (Report) tổng quan về thái độ và năng lực.
      \item Gửi đánh giá cho học viên.
    \end{enumerate}
  },
  dongphu = {Không.},
  haudk = {Báo cáo đánh giá cá nhân được lưu vào hồ sơ học viên.}
}

% --- UC035 ---
\UseCaseBox{
  ma = {UC035},
  ten = {Theo dõi tiến độ học tập},
  tacnhan = {Giáo viên},
  muctieu = {Nhận diện tình hình chung của lớp, phát hiện học viên tụt hậu.},
  mota = {Xem tỷ lệ nộp bài, biểu đồ điểm trung bình của cả lớp.},
  yeucau = {Không.},
  tiendk = {Đăng nhập Teacher.},
  dongchinh = {
    \begin{enumerate}[leftmargin=*, topsep=0pt, itemsep=0pt]
      \item Truy cập "Báo cáo lớp học".
      \item Hệ thống hiển thị Dashboard: Tỷ lệ hoàn thành bài tập (\%), tỷ lệ học viên dưới trung bình.
      \item Giáo viên dựa vào đó để điều chỉnh phương pháp dạy.
    \end{enumerate}
  },
  dongphu = {Không.},
  haudk = {Giáo viên có dữ liệu để quản trị chất lượng lớp học.}
}

% --- UC036 ---
\UseCaseBox{
  ma = {UC036},
  ten = {Xem báo cáo và thống kê học tập},
  tacnhan = {Quản lý (Admin), Giáo viên},
  muctieu = {Quản lý bức tranh tổng thể hoạt động của toàn trung tâm/nhiều lớp.},
  mota = {Xuất báo cáo hiệu suất, thống kê số lượng bài được giao, điểm số trung bình các khóa.},
  yeucau = {Không.},
  tiendk = {Đăng nhập Admin hoặc Teacher.},
  dongchinh = {
    \begin{enumerate}[leftmargin=*, topsep=0pt, itemsep=0pt]
      \item Quản lý truy cập "Thống kê báo cáo".
      \item Thiết lập bộ lọc (Từ tháng X đến tháng Y, theo chi nhánh/lớp/giáo viên).
      \item Hệ thống xuất ra các báo cáo dạng số liệu và biểu đồ (KPI).
      \item Nhấn "Xuất PDF/Excel" để tải báo cáo.
    \end{enumerate}
  },
  dongphu = {
    \begin{itemize}[leftmargin=*, topsep=0pt, itemsep=0pt]
      \item Không có dữ liệu trong khoảng thời gian đã chọn $\rightarrow$ Hiển thị "Không có dữ liệu".
    \end{itemize}
  },
  haudk = {Tải về được file báo cáo thống kê chính xác phục vụ cho việc vận hành trung tâm.}
}

\end{document}